\section{Summary}

This template has demonstrated how a dissertation can be divided into modular LaTeX files while retaining a unified structure in \texttt{thesis\_main.tex}. The introduction established a problem and previewed the document. The literature review illustrated thematic synthesis and reference management. The three example study chapters demonstrated mathematical notation, tables, standalone and inline TikZ graphics, an algorithm, theorem environments, formal proofs, and a PGFPlots chart.

\section{Connecting Findings to Contributions}

An actual conclusion should answer the research questions rather than merely repeat chapter headings. It should explain the overall contribution, synthesize how the individual studies support that contribution, and distinguish findings supported by evidence from broader implications. References such as Equation~\eqref{eq:weighted-score}, Algorithm~\ref{alg:select-best}, and Figure~\ref{fig:performance-plot} demonstrate that labels remain available across chapter files.

\section{Limitations}

The content in this project is instructional. Numerical values are fictional, the research questions are generic, and the mathematical model is deliberately simple. The included \texttt{mstogs.cls} class reflects a particular institutional formatting system and may not satisfy another graduate school's current requirements. Before using the template for submission, the writer must compare it with the latest official formatting guide and replace or revise the class file if necessary.

\section{Future Research Directions}

Future-work statements should emerge from the findings and limitations. They may propose validation in a new context, collection of additional evidence, improvement of a method, comparison with stronger alternatives, or investigation of a newly observed question. Each direction should explain why it matters and how it extends the completed dissertation.

For this template project, useful extensions include adding discipline-specific chapter structures, automating glossary generation, integrating source-code listings, adding an acronyms package, and configuring continuous integration to compile the dissertation whenever changes are pushed to a repository.

\section{Final Checklist}

Before submission, confirm that all placeholders have been replaced; every citation resolves; figures and tables are readable; cross-references point to the correct elements; permissions and attribution requirements are satisfied; front matter follows institutional policy; and the final PDF has been inspected page by page.
